\IfFileExists{stacks-project.cls}{%
\documentclass{stacks-project}
}{%
\documentclass{amsart}
}
\usepackage{amssymb}

% For dealing with references we use the comment environment
\usepackage{verbatim}
\usepackage{fancyvrb}
\newenvironment{reference}{\comment}{\endcomment}
%\newenvironment{reference}{}{}
\newenvironment{slogan}{\comment}{\endcomment}
\newenvironment{history}{\comment}{\endcomment}

% For source figures retained as assets
\usepackage{graphicx}

% For commutative diagrams we use Xy-pic
\usepackage[all]{xy}

% We use 2cell for 2-commutative diagrams.
\xyoption{2cell}
\UseAllTwocells

% We use multicol for the list of chapters between chapters
\usepackage{multicol}

% This is generally recommended for better output
\usepackage{lmodern}
\usepackage[T1]{fontenc}
\usepackage{textalpha}

% For cross-file-references
% We prefer the snapshot tmp/xref/#2.aux,
% which latexmkrc replaces atomically, so
% that documents compiled simultaneously
% never read a half-written .aux file.
% The last argument keeps links to #2.pdf.
\usepackage{xr-hyper}
\newcommand{\maybeexternaldocument}[2]{%
  \IfFileExists{#2.aux}{%
    \externaldocument[#1]{#2}%
  }{%
    \PackageWarning{stack}{Missing #2.aux; external references with prefix #1 unavailable}%
  }%
}
\let\externaldocumentorig\externaldocument
\renewcommand{\externaldocument}[2][]{%
  \IfFileExists{tmp/xref/#2.aux}{%
    \externaldocumentorig[#1]{tmp/xref/#2}[#2.pdf]%
  }{%
  \IfFileExists{#2.aux}{%
    \externaldocumentorig[#1]{#2}%
  }{%
    \PackageWarning{stack}{Missing #2.aux; external references with prefix #1 unavailable}%
  }%
  }%
}

% Package for hypertext links:
\usepackage{hyperref}

% For any local file, say "hello.tex" you want to link to please
% use \externaldocument[hello-]{hello}
\externaldocument[introduction-]{introduction}
\externaldocument[conventions-]{conventions}
\externaldocument[sets-]{sets}
\externaldocument[categories-]{categories}
\externaldocument[topology-]{topology}
\externaldocument[sheaves-]{sheaves}
\externaldocument[sites-]{sites}
\externaldocument[stacks-]{stacks}
\externaldocument[fields-]{fields}
\externaldocument[algebra-]{algebra}
\externaldocument[brauer-]{brauer}
\externaldocument[homology-]{homology}
\externaldocument[derived-]{derived}
\externaldocument[simplicial-]{simplicial}
\externaldocument[more-algebra-]{more-algebra}
\externaldocument[smoothing-]{smoothing}
\externaldocument[modules-]{modules}
\externaldocument[sites-modules-]{sites-modules}
\externaldocument[injectives-]{injectives}
\externaldocument[cohomology-]{cohomology}
\externaldocument[sites-cohomology-]{sites-cohomology}
\externaldocument[dga-]{dga}
\externaldocument[dpa-]{dpa}
\externaldocument[sdga-]{sdga}
\externaldocument[hypercovering-]{hypercovering}
\externaldocument[schemes-]{schemes}
\externaldocument[constructions-]{constructions}
\externaldocument[properties-]{properties}
\externaldocument[morphisms-]{morphisms}
\externaldocument[coherent-]{coherent}
\externaldocument[divisors-]{divisors}
\externaldocument[limits-]{limits}
\externaldocument[varieties-]{varieties}
\externaldocument[topologies-]{topologies}
\externaldocument[descent-]{descent}
\externaldocument[perfect-]{perfect}
\externaldocument[more-morphisms-]{more-morphisms}
\externaldocument[flat-]{flat}
\externaldocument[groupoids-]{groupoids}
\externaldocument[more-groupoids-]{more-groupoids}
\externaldocument[etale-]{etale}
\externaldocument[chow-]{chow}
\externaldocument[intersection-]{intersection}
\externaldocument[pic-]{pic}
\externaldocument[weil-]{weil}
\externaldocument[adequate-]{adequate}
\externaldocument[dualizing-]{dualizing}
\externaldocument[duality-]{duality}
\externaldocument[discriminant-]{discriminant}
\externaldocument[derham-]{derham}
\externaldocument[local-cohomology-]{local-cohomology}
\externaldocument[algebraization-]{algebraization}
\externaldocument[curves-]{curves}
\externaldocument[resolve-]{resolve}
\externaldocument[models-]{models}
\externaldocument[functors-]{functors}
\externaldocument[equiv-]{equiv}
\externaldocument[pione-]{pione}
\externaldocument[etale-cohomology-]{etale-cohomology}
\externaldocument[proetale-]{proetale}
\externaldocument[relative-cycles-]{relative-cycles}
\externaldocument[more-etale-]{more-etale}
\externaldocument[trace-]{trace}
\externaldocument[crystalline-]{crystalline}
\externaldocument[spaces-]{spaces}
\externaldocument[spaces-properties-]{spaces-properties}
\externaldocument[spaces-morphisms-]{spaces-morphisms}
\externaldocument[decent-spaces-]{decent-spaces}
\externaldocument[spaces-cohomology-]{spaces-cohomology}
\externaldocument[spaces-limits-]{spaces-limits}
\externaldocument[spaces-divisors-]{spaces-divisors}
\externaldocument[spaces-over-fields-]{spaces-over-fields}
\externaldocument[spaces-topologies-]{spaces-topologies}
\externaldocument[spaces-descent-]{spaces-descent}
\externaldocument[spaces-perfect-]{spaces-perfect}
\externaldocument[spaces-more-morphisms-]{spaces-more-morphisms}
\externaldocument[spaces-flat-]{spaces-flat}
\externaldocument[spaces-groupoids-]{spaces-groupoids}
\externaldocument[spaces-more-groupoids-]{spaces-more-groupoids}
\externaldocument[bootstrap-]{bootstrap}
\externaldocument[spaces-pushouts-]{spaces-pushouts}
\externaldocument[spaces-chow-]{spaces-chow}
\externaldocument[groupoids-quotients-]{groupoids-quotients}
\externaldocument[spaces-more-cohomology-]{spaces-more-cohomology}
\externaldocument[spaces-simplicial-]{spaces-simplicial}
\externaldocument[spaces-duality-]{spaces-duality}
\externaldocument[formal-spaces-]{formal-spaces}
\externaldocument[restricted-]{restricted}
\externaldocument[spaces-resolve-]{spaces-resolve}
\externaldocument[formal-defos-]{formal-defos}
\externaldocument[defos-]{defos}
\externaldocument[cotangent-]{cotangent}
\externaldocument[examples-defos-]{examples-defos}
\externaldocument[algebraic-]{algebraic}
\externaldocument[examples-stacks-]{examples-stacks}
\externaldocument[stacks-sheaves-]{stacks-sheaves}
\externaldocument[criteria-]{criteria}
\externaldocument[artin-]{artin}
\externaldocument[quot-]{quot}
\externaldocument[stacks-properties-]{stacks-properties}
\externaldocument[stacks-morphisms-]{stacks-morphisms}
\externaldocument[stacks-limits-]{stacks-limits}
\externaldocument[stacks-cohomology-]{stacks-cohomology}
\externaldocument[stacks-perfect-]{stacks-perfect}
\externaldocument[stacks-introduction-]{stacks-introduction}
\externaldocument[stacks-more-morphisms-]{stacks-more-morphisms}
\externaldocument[stacks-geometry-]{stacks-geometry}
\externaldocument[moduli-]{moduli}
\externaldocument[moduli-curves-]{moduli-curves}
\externaldocument[examples-]{examples}
\externaldocument[exercises-]{exercises}
\externaldocument[guide-]{guide}
\externaldocument[desirables-]{desirables}
\externaldocument[coding-]{coding}
\externaldocument[obsolete-]{obsolete}
\externaldocument[fdl-]{fdl}
\externaldocument[index-]{index}

% Heap Project documents
\externaldocument[k3-]{k3}
\externaldocument[torelli-]{torelli}
\externaldocument[linear-algebra-]{linear-algebra}
\externaldocument[groups-]{groups}
\externaldocument[complex-analysis-]{complex-analysis}
\externaldocument[functional-analysis-]{functional-analysis}
\externaldocument[categories-]{categories}
\externaldocument[commutative-algebra-]{commutative-algebra}
\externaldocument[ringed-spaces-]{ringed-spaces}
\externaldocument[homological-algebra-]{homological-algebra}
\externaldocument[lie-algebras-]{lie-algebras}
\externaldocument[representation-theory-]{representation-theory}
\externaldocument[smooth-manifolds-]{smooth-manifolds}
\externaldocument[differential-topology-]{differential-topology}
\externaldocument[algebraic-topology-]{algebraic-topology}
\externaldocument[differential-geometry-]{differential-geometry}
\externaldocument[algebraic-geometry-]{algebraic-geometry}
\externaldocument[symplectic-geometry-]{symplectic-geometry}
\externaldocument[complex-geometry-]{complex-geometry}
\externaldocument[hodge-theory-]{hodge-theory}
\externaldocument[%
variations-of-hodge-structures-]{%
variations-of-hodge-structures}
\externaldocument[%
homological-mirror-symmetry-]{%
homological-mirror-symmetry}
\maybeexternaldocument{deformations-}{deformations}
\externaldocument[vertex-algebras-]{vertex-algebras}
\externaldocument[springer-]{springer}
\externaldocument[infty-categories-]{infty-categories}
\externaldocument[seminars-differential-geometry-]{%
seminars-differential-geometry}
\externaldocument[seminars-complex-geometry-]{%
seminars-complex-geometry}
\externaldocument[seminars-algebraic-geometry-]{%
seminars-algebraic-geometry}
\externaldocument[seminars-others-]{%
seminars-others}
\externaldocument[%
geometric-prequantization-]{%
geometric-prequantization}
\externaldocument[surfaces-]{surfaces}
\externaldocument[birational-maps-commutative-algebra-]{birational-maps-commutative-algebra}
\externaldocument[cremona-transformations-]{cremona-transformations}
\externaldocument[derived-categories-of-sheaves-]{derived-categories-of-sheaves}
\externaldocument[geometric-invariant-theory-]{geometric-invariant-theory}
\externaldocument[geometric-stability-conditions-and-group-actions-]{geometric-stability-conditions-and-group-actions}
\externaldocument[moduli-spaces-of-sheaves-]{moduli-spaces-of-sheaves}
\externaldocument[bridgeland-stability-general-theory-]{bridgeland-stability-general-theory}
\externaldocument[hilbert-schemes-]{hilbert-schemes}
\externaldocument[noncommutative-abelian-surfaces-and-kummer-type-hyperkahler-varieties-]{noncommutative-abelian-surfaces-and-kummer-type-hyperkahler-varieties}
\externaldocument[mumford-tate-groups-in-hodge-theory-]{mumford-tate-groups-in-hodge-theory}
\externaldocument[hodge-birational-atoms-2026-]{hodge-birational-atoms-2026}
\externaldocument[xxii-egd-2026-]{%
xxii-egd-2026}
\externaldocument[pde-]{pde}
\externaldocument[probability-]{probability}

\externaldocument[linguistics-]{linguistics}
\externaldocument[genealogia-familiar-]{%
genealogia-familiar}

\externaldocument[%
16th-alga-meeting-2026-]{%
16th-alga-meeting-2026}

\externaldocument[%
iv-jornadas-lefschetz-2026-]{%
iv-jornadas-lefschetz-2026}

\externaldocument[reading-the-masters-]{%
reading-the-masters}

\externaldocument[real-analysis-]{%
real-analysis}

\externaldocument[optimal-transport-]{%
optimal-transport}

\externaldocument[history-]{history}

\externaldocument[seiberg-witten-]{seiberg-witten}

\externaldocument[history-philosophy-biology-]{history-philosophy-biology}

\externaldocument[chemistry-]{chemistry}

\externaldocument[molecular-cell-biology-]{molecular-cell-biology}

\externaldocument[microbiology-]{microbiology}

\externaldocument[protistology-]{protistology}

\externaldocument[mycology-]{mycology}

\externaldocument[botany-]{botany}

\externaldocument[zoology-]{zoology}

\externaldocument[animal-histology-]{animal-histology}

\externaldocument[animal-physiology-]{animal-physiology}

\externaldocument[systematics-biogeography-]{systematics-biogeography}

\externaldocument[ecology-natural-resources-]{ecology-natural-resources}

\externaldocument[earth-sciences-]{earth-sciences}

\externaldocument[biostatistics-]{biostatistics}

\externaldocument[mathematical-systems-biology-]{mathematical-systems-biology}

\externaldocument[cimpa-el-salvador-2024-]{cimpa-el-salvador-2024}

% Theorem environments.
%
\theoremstyle{plain}
\newtheorem{theorem}[subsection]{Theorem}
\newtheorem{proposition}[subsection]{Proposition}
\newtheorem{lemma}[subsection]{Lemma}
\newtheorem{corollary}[subsection]{Corollary}

\theoremstyle{definition}
\newtheorem{definition}[subsection]{Definition}
\newtheorem{example}[subsection]{Example}
\newtheorem{exercise}[subsection]{Exercise}
\newtheorem{situation}[subsection]{Situation}

\theoremstyle{remark}
\newtheorem{remark}[subsection]{Remark}
\newtheorem{remarks}[subsection]{Remarks}

\numberwithin{equation}{subsection}

% Macros
%
\def\lim{\mathop{\mathrm{lim}}\nolimits}
\def\colim{\mathop{\mathrm{colim}}\nolimits}
\def\Spec{\mathop{\mathrm{Spec}}}
\def\Hom{\mathop{\mathrm{Hom}}\nolimits}
\def\Ext{\mathop{\mathrm{Ext}}\nolimits}
\def\SheafHom{\mathop{\mathcal{H}\!\mathit{om}}\nolimits}
\def\SheafExt{\mathop{\mathcal{E}\!\mathit{xt}}\nolimits}
\def\Sch{\mathit{Sch}}
\def\Mor{\mathop{\mathrm{Mor}}\nolimits}
\def\Ob{\mathop{\mathrm{Ob}}\nolimits}
\def\Sh{\mathop{\mathit{Sh}}\nolimits}
\def\NL{\mathop{N\!L}\nolimits}
\def\CH{\mathop{\mathrm{CH}}\nolimits}
\def\proetale{{pro\text{-}\acute{e}tale}}
\def\etale{{\acute{e}tale}}
\def\QCoh{\mathit{QCoh}}
\def\Ker{\mathop{\mathrm{Ker}}}
\def\Im{\mathop{\mathrm{Im}}}
\def\Coker{\mathop{\mathrm{Coker}}}
\def\Coim{\mathop{\mathrm{Coim}}}

% Boxtimes
%
\DeclareMathSymbol{\boxtimes}{\mathbin}{AMSa}{"02}

%
% Macros for moduli stacks/spaces
%
\def\QCohstack{\mathcal{QC}\!\mathit{oh}}
\def\Cohstack{\mathcal{C}\!\mathit{oh}}
\def\Spacesstack{\mathcal{S}\!\mathit{paces}}
\def\Quotfunctor{\mathrm{Quot}}
\def\Hilbfunctor{\mathrm{Hilb}}
\def\Curvesstack{\mathcal{C}\!\mathit{urves}}
\def\Polarizedstack{\mathcal{P}\!\mathit{olarized}}
\def\Complexesstack{\mathcal{C}\!\mathit{omplexes}}
% \Pic is the operator that assigns to X its picard group, usage \Pic(X)
% \Picardstack_{X/B} denotes the Picard stack of X over B
% \Picardfunctor_{X/B} denotes the Picard functor of X over B
\def\Pic{\mathop{\mathrm{Pic}}\nolimits}
\def\Picardstack{\mathcal{P}\!\mathit{ic}}
\def\Picardfunctor{\mathrm{Pic}}
\def\Deformationcategory{\mathcal{D}\!\mathit{ef}}


\begin{document}

\title{Microbiology}
\maketitle

\phantomsection
\label{section-phantom}

\tableofcontents

\section{What this chapter covers}
\label{section-overview}

These notes cover the biology of Bacteria and
Archaea: the prokaryotic cell and how it
differs from the eukaryotic cell, the cell
envelope and the Gram stain, motility and the
prokaryotic cytoskeleton, inclusions and
endospores, metabolic diversity and the
Archaea, gene regulation and horizontal gene
transfer, growth and culture in the
laboratory, and a short list of genera that
matter in medicine, agriculture and the
environment.

\textbf{Sources of these notes.} The course
was
\emph{Biolog\'ia de Procariontes}, first
semester of the Biology degree at the
Facultad de Ciencias, UNAM, in the autumn of
2013 (UNAM semester 2014-1). What survives is
the course's laboratory manual with nine
practicals, the lecture slides on cell
structure, the cytoskeleton and motility, and
on gene regulation and gene transfer,
students' slides on metabolism and Archaea,
and Daniel's own work: methods write-ups for
the practicals, short homework on the Gram
stain and on selective media, a take-home
exam (with an essay on the term
``prokaryote''), a team project on
\emph{Vibrio} and \emph{Klebsiella}, a study
list of genera for the final exam and an
essay on pollution and biogeochemical cycles.
The notes were checked and completed with
Brock's textbook \cite{madigan2012brock},
Tortora's \cite{tortora2013microbiology},
Prescott's \cite{willey2008prescott} and the
2009 exchange between Pace and Whitman
\cite{pace2009problems, whitman2009modern,
pace2009rebuttal, whitman2009rebuttal}.

\section{Prokaryotic and eukaryotic cells}
\label{section-prokaryotic-eukaryotic}

The prokaryotic cell \cite[Section
2.5]{madigan2012brock}
\cite{whitman2009modern}:
\begin{itemize}
\item Its DNA lies in the cytoplasm, in a
region called the nucleoid, not inside a
nucleus. Because of this, transcription and
translation are coupled: ribosomes start
translating an mRNA while it is still being
made.
\item It is usually small: a typical rod is
1--5 $\mu$m long and about 1 $\mu$m wide.
There are exceptions, such as the sulfur
bacterium \emph{Thiomargarita}, up to 750
$\mu$m across.
\item Its cytoplasmic membrane is
multifunctional: respiration and
photosynthesis take place on it and generate
the proton motive force that drives ATP
synthesis, transport and the flagella.
\item It divides by binary fission, not by
mitosis; its genome is usually a single
circular chromosome, often with plasmids.
\item It does not take up particles by
phagocytosis.
\end{itemize}
The eukaryotic cell has a nucleus with a
double envelope, divides by mitosis, has
membrane-bound organelles (mitochondria and,
in algae and plants, chloroplasts), an actin
and tubulin cytoskeleton, and flagella with
the 9+2 arrangement of microtubules that beat
by dynein motors \cite[Chapter
1]{lee2008phycology}. Prokaryotic flagella
are completely different: they are rigid
helical filaments that rotate
(Section~\ref{section-motility}).

Comparing the sequences of small-subunit
(16S) ribosomal RNA, Woese showed in 1977
that living things fall into three domains:
Bacteria, Archaea and Eukarya \cite[Section
2.7]{madigan2012brock}
\cite{pace2009problems}. Bacteria and Archaea
are both prokaryotic in cell structure, but
they are not each other's closest relatives:
in the rRNA tree the Archaea are a sister
group of the Eukarya.

\section{The debate on the term
``prokaryote''}
\label{section-prokaryote-debate}

In April 2009 the Journal of Bacteriology
published four short papers in which Pace and
Whitman argued over the word ``prokaryote'';
the course set them as a take-home exam
question.

Pace's case \cite{pace2009problems,
pace2009rebuttal}. Stanier and van Niel
(1962) defined prokaryotes by the absence of
a nucleus, and Stanier's textbook (1963)
turned this into an evolutionary model:
prokaryotes are a natural group from which
eukaryotes arose. Molecular phylogeny refuted
this model. There are three domains; Archaea
and Eukarya share a common ancestor that
excludes Bacteria, so ``prokaryotes'' are not
a group of related organisms, and no group of
living organisms preceded the eukaryotes.
Moreover, the definition is negative: it says
only what a prokaryote is not. Pace concluded
that the word should be retired.

Whitman's case \cite{whitman2009modern,
whitman2009rebuttal}. Stanier and van Niel
defined the prokaryote by cell structure, not
by phylogeny: no internal membranes around
the nuclear material or the respiratory and
photosynthetic machinery, division by fission
instead of mitosis, and a peptidoglycan wall.
The modern concept has positive content:
coupled transcription and translation, small
cells with a fast diffusion of molecules, a
multifunctional cytoplasmic membrane, no
phagocytosis. So ``prokaryote'' names a real
type of cellular organization, shared by
Bacteria and Archaea, whatever their
phylogeny. Brock's textbook takes the same
line, and adds that ``prokaryote'' should
carry no evolutionary connotation and does
not mean ``primitive'' \cite[Section
2.5]{madigan2012brock}.

\begin{remark}
\label{remark-daniel-position}
In his take-home exam Daniel took a middle
position. He agreed with Pace that a negative
definition is logically weak and that the
term must not suggest that prokaryotes are
ancestors of eukaryotes. But the absence of a
nucleus can be restated positively, as
nucleic acids lying free in the cytoplasm and
as transcription, translation and replication
happening at the same time and place; and
Bacteria and Archaea share such traits, which
eukaryotes lack. He concluded that
``prokaryote'' is a fair name for Bacteria
and Archaea taken together, as a type of cell
and not a branch of the tree of life, to be
used and taught with care.
\end{remark}

\section{Cell envelope and the Gram stain}
\label{section-gram-stain}

The bacterial cell wall is made of
peptidoglycan, a mesh of glycan chains
(alternating N-acetylglucosamine and
N-acetylmuramic acid) cross-linked by short
peptides \cite[Sections
3.6--3.8]{madigan2012brock}.
\begin{itemize}
\item Gram-positive bacteria have a thick
wall, up to 90\% peptidoglycan, with teichoic
and lipoteichoic acids embedded in it.
\item Gram-negative bacteria have a thin
peptidoglycan layer (about 10\% of the wall)
and, outside it, an outer membrane: a second
lipid bilayer whose outer leaflet contains
lipopolysaccharide (LPS, whose lipid A is the
endotoxin). Porins in the outer membrane let
small molecules through. Between the two
membranes lies the periplasm.
\item Archaea have no peptidoglycan. Some
methanogens have a wall of pseudomurein; many
Archaea have only an S-layer, a
paracrystalline layer of protein or
glycoprotein. Their membrane lipids are
ether-linked isoprenoids, not the
ester-linked fatty acids of Bacteria and
Eukarya.
\end{itemize}

The Gram stain, done in the course's second
practical, is the basic differential stain of
bacteriology \cite[Section
2.2]{madigan2012brock}:
\begin{enumerate}
\item stain a heat-fixed smear with crystal
violet (1 minute): all cells are purple;
\item add iodine (Lugol's solution), which
forms an insoluble crystal violet--iodine
complex inside the cells;
\item decolorize with alcohol (in the
course's manual, alcohol--acetone);
\item counterstain with safranin.
\end{enumerate}
Gram-positive cells stay purple-violet;
Gram-negative cells lose the violet and
appear pink to red. The reason is the wall
\cite[Section 3.7]{madigan2012brock}: alcohol
dehydrates the thick peptidoglycan of
Gram-positive cells and closes its pores,
trapping the complex, while in Gram-negative
cells it dissolves into the lipid-rich outer
membrane and washes the complex out of the
thin wall.

\section{Flagella, motility and the
cytoskeleton}
\label{section-motility}

The bacterial flagellum \cite[Section
3.13]{madigan2012brock} has three parts: a
helical filament made of the protein
flagellin, a hook, and a basal body anchored
in the envelope that works as a rotary motor.
Its stator is formed by the Mot proteins, and
the energy for rotation comes from protons
flowing into the cell through them, that is,
from the proton motive force. The flagellum
pushes or pulls the cell by rotating.
Arrangements:
\begin{itemize}
\item polar: one flagellum at one end
(monotrichous), or a tuft at one end
(lophotrichous) or at both ends
(amphitrichous);
\item peritrichous: flagella all over the
surface;
\item endoflagella of spirochetes: flagella
that stay in the periplasm and turn the whole
cell like a corkscrew \cite[Chapter
18]{madigan2012brock}.
\end{itemize}
Archaeal flagella also rotate, but they are
about half as thick and are built from
several flagellins unrelated to the bacterial
one. Other ways of moving are gliding (for
example \emph{Flavobacterium}) and twitching
by extension and retraction of type IV pili
\cite[Sections 3.13--3.14]{madigan2012brock}.
Motile cells move towards attractants and
away from repellents (taxis): chemotaxis,
phototaxis, and magnetotaxis in bacteria with
magnetosomes \cite[Section
3.15]{madigan2012brock}.

Bacteria have a cytoskeleton made of
homologues of the eukaryotic proteins
\cite[Sections 5.2--5.3]{madigan2012brock}
\cite{wickstead2011evolution}:
\begin{itemize}
\item FtsZ, related to tubulin, forms a ring
at the division plane (the Z ring) that
organizes the division machinery. It occurs
in Bacteria and Archaea, and FtsZ-type
proteins are also found in chloroplasts and
mitochondria, a trace of their bacterial
origin. The Min proteins oscillate from pole
to pole and keep the Z ring in the middle of
the cell \cite[Chapter
6]{willey2008prescott}.
\item MreB, related to actin, forms filaments
under the membrane and directs wall synthesis
so that rod-shaped cells grow along their
long axis.
\item Crescentin, related to the intermediate
filaments (keratins) of animals, lies on the
concave side of the curved cells of
\emph{Caulobacter crescentus} and gives them
their shape.
\item Other actin homologues include ParM,
which segregates plasmids at division, and
MamK, which arranges magnetosomes in a chain.
\end{itemize}

\section{Inclusions, microcompartments and
endospores}
\label{section-inclusions}

The cytoplasm holds several kinds of
inclusions \cite[Sections
3.10--3.11]{madigan2012brock}:
\begin{itemize}
\item storage polymers:
poly-$\beta$-hydroxybutyrate (PHB) and other
polyhydroxyalkanoates for carbon and energy,
glycogen, polyphosphate granules, and
globules of elemental sulfur in
sulfur-oxidizing bacteria;
\item magnetosomes: membrane-bound crystals
of magnetite ($\mathrm{Fe_3O_4}$), arranged
in chains, that orient magnetotactic bacteria
in the Earth's magnetic field;
\item gas vesicles: spindle-shaped hollow
structures made only of protein, about 300 to
more than 1000 nm long and 45--120 nm wide,
permeable to gases but not to water. They
lower the density of the cell and let
planktonic Bacteria and Archaea float at a
chosen depth; cyanobacteria such as
\emph{Anabaena} and \emph{Microcystis} with
gas vesicles form blooms on lakes. They have
never been found in eukaryotes.
\end{itemize}
Carboxysomes are polyhedral bodies about 100
nm across with a thin protein shell, found in
many autotrophs that fix $\mathrm{CO_2}$ by
the Calvin cycle \cite[Section
13.12]{madigan2012brock}. They pack about 250
molecules of RubisCO together with carbonic
anhydrase, which turns bicarbonate into
$\mathrm{CO_2}$; the shell keeps the
$\mathrm{CO_2}$ in and limits access of
$\mathrm{O_2}$, so RubisCO carboxylates
instead of oxygenating. They are the best
known example of a protein-shelled
microcompartment.

Endospores \cite[Section
3.12]{madigan2012brock} are dormant cells
formed inside a mother cell by
\emph{Bacillus}, \emph{Clostridium} and
relatives when nutrients run out. From
outside in, an endospore has an exosporium,
protein spore coats, a cortex of loosely
cross-linked peptidoglycan, and a dehydrated
core rich in calcium bound to dipicolinic
acid (about 10\% of the spore's dry weight)
and in small acid-soluble proteins that
protect the DNA. Endospores survive boiling,
drying, radiation and disinfectants for very
long times and germinate when conditions
improve. Their heat resistance is why
laboratory sterilization uses an autoclave,
steam under pressure at about 121 $^\circ$C,
rather than boiling.

\section{Metabolic diversity and the Archaea}
\label{section-metabolism-archaea}

By energy source, organisms are \cite[Section
2.8]{madigan2012brock}:
\begin{itemize}
\item chemoorganotrophs, which oxidize
organic compounds (most cultured microbes);
\item chemolithotrophs, which oxidize
inorganic compounds such as $\mathrm{H_2}$,
$\mathrm{H_2S}$, $\mathrm{NH_3}$ or
$\mathrm{Fe^{2+}}$; this occurs only in
prokaryotes;
\item phototrophs: oxygenic (cyanobacteria,
like algae and plants) or anoxygenic (purple
and green bacteria, heliobacteria).
\end{itemize}
By carbon source they are heterotrophs
(organic carbon) or autotrophs
($\mathrm{CO_2}$). Sugars are broken down by
glycolysis (the Embden--Meyerhof pathway) or,
in \emph{Pseudomonas} and \emph{Zymomonas},
by the Entner--Doudoroff pathway
\cite[Section 14.2]{madigan2012brock}.
Without an external electron acceptor, cells
ferment, producing for example ethanol,
lactate, butanol and acetone
(\emph{Clostridium}) or propionate
(\emph{Propionibacterium}); with one, they
respire, aerobically with $\mathrm{O_2}$ or
anaerobically with nitrate, sulfate, iron and
other acceptors.

The Archaea \cite[Chapter
19]{madigan2012brock} have ether-linked
lipids, no peptidoglycan, and transcription
and translation machinery closer to that of
eukaryotes. Many are extremophiles. The two
main phyla in the course's sources were:
\begin{itemize}
\item Euryarchaeota: extreme halophiles such
as \emph{Halobacterium salinarum}, which
grows in near-saturated brine and makes ATP
with light through bacteriorhodopsin, a
light-driven proton pump that colours the
cells purple; the methanogens; and
hyperthermophiles such as \emph{Pyrococcus}
(growth from 70 to 106 $^\circ$C, optimum 100
$^\circ$C) and the methanogen
\emph{Methanopyrus kandleri}, which grows up
to 122 $^\circ$C. \emph{Nanoarchaeum
equitans}, one of the smallest cells known
(genome 0.49 Mb), lives attached to the
archaeon \emph{Ignicoccus}.
\item Crenarchaeota: mostly hyperthermophiles
of hot springs and submarine vents, such as
\emph{Sulfolobus}, but also abundant marine
and soil Archaea that oxidize ammonia.
\end{itemize}
Methanogens (for example
\emph{Methanosarcina}) are strict anaerobes
that produce methane, for example from
hydrogen and carbon dioxide,
$$
\mathrm{CO_2} + 4\,\mathrm{H_2} \rightarrow
\mathrm{CH_4} + 2\,\mathrm{H_2O},
$$
and use unusual coenzymes such as the
nickel-containing coenzyme
$\mathrm{F_{430}}$.

\section{Gene regulation and horizontal gene
transfer}
\label{section-genetics}

\begin{definition}
\label{definition-operon}
An \emph{operon} is a group of genes
transcribed together from one promoter into a
single mRNA, under the control of a common
regulatory site, the \emph{operator}.
\end{definition}

Bacteria adjust their enzymes to the food
available mainly by controlling transcription
\cite[Sections 8.3--8.5 and
8.16]{madigan2012brock}:
\begin{itemize}
\item Induction. In the \emph{lac} operon of
\emph{E. coli}, which encodes the uptake and
cleavage of lactose, the Lac repressor binds
the operator and blocks transcription. When
lactose is present, its isomer allolactose
(the inducer) binds the repressor, which
releases the DNA.
\item Catabolite repression. Even with
lactose present, the \emph{lac} operon is
barely expressed while glucose is available:
glucose keeps the level of cyclic AMP low,
and the activator protein CRP needs cyclic
AMP to bind the promoter region. Glucose is
thus not a corepressor; it acts through the
lack of an activator.
\item Repression. In biosynthetic operons the
end product acts as a corepressor: it binds
the repressor and lets it bind the operator
(arginine for the arginine genes, tryptophan
for the \emph{trp} operon).
\item Attenuation. In the \emph{trp} operon,
transcription also stops early in a leader
region when there is enough tryptophan,
because the speed at which ribosomes
translate a short leader peptide decides
which RNA hairpin forms. This works only
because transcription and translation are
coupled.
\end{itemize}

Prokaryotes reproduce asexually, but they
exchange genes between cells, even across
species, by horizontal gene transfer. The
course presented this as the main source of
their genetic novelty. There are three
mechanisms \cite[Sections
10.7--10.9]{madigan2012brock}:
\begin{itemize}
\item Transformation: uptake of free DNA from
the environment by cells in a physiological
state called competence. Griffith discovered
it in the 1920s with pneumococcus;
\emph{Bacillus}, \emph{Streptococcus},
\emph{Haemophilus} and \emph{Neisseria} are
naturally competent.
\item Transduction: transfer of bacterial DNA
by a bacteriophage. In generalized
transduction any part of the donor's genome
can be packed by mistake into a phage head;
in specialized transduction a temperate phage
(such as $\lambda$ in \emph{E. coli}) carries
only the genes next to its integration site.
\item Conjugation: transfer of DNA through
cell contact, encoded by plasmids such as the
F plasmid, whose \emph{tra} genes make the
sex pilus and the transfer machinery;
transfer starts at the origin of transfer
(\emph{oriT}). When F is integrated into the
chromosome (Hfr strains), chromosomal genes
are transferred too.
\end{itemize}
Plasmids replicate independently of the
chromosome and often carry genes for
antibiotic resistance, toxins or unusual
metabolism, which is how such traits spread
through bacterial populations.

\section{Growth and culture}
\label{section-growth}

Bacteria grow by binary fission, so a
population doubles once per generation time
$g$ (about 20 minutes for \emph{E. coli} in
rich medium, according to the course's
manual). After $n = t/g$ generations a
population of $N_0$ cells has
$$
N = N_0 \, 2^{n}
$$
cells \cite[Sections
5.5--5.6]{madigan2012brock}. In a closed
culture the growth curve has four phases
\cite[Section 5.7]{madigan2012brock}: lag
(cells adjust to the medium), exponential
(constant doubling), stationary (nutrients
run out or wastes accumulate; no net growth)
and death.

Culture media are defined (known chemical
composition) or complex (digests of protein,
yeast extract), and may be enriched for
fastidious organisms. A \emph{selective
medium} contains compounds that inhibit some
organisms but not others; a
\emph{differential medium} contains an
indicator that shows whether a given reaction
took place \cite[Section
4.2]{madigan2012brock}. The course used:
\begin{itemize}
\item MacConkey agar: bile salts and crystal
violet inhibit Gram-positive bacteria
(selective); lactose with the pH indicator
neutral red marks lactose fermenters
(differential);
\item EMB (eosin--methylene blue) agar, for
enteric bacteria;
\item TCBS agar: thiosulfate, citrate, bile
and a high pH favour \emph{Vibrio};
\item minimal media with a single carbon
source (citrate, lactose, arabinose, sucrose,
rhamnose, dulcitol, salicin, raffinose) and
TSI (triple sugar iron) tubes, to define
biotypes.
\end{itemize}

The manual's nine practicals were: laboratory
equipment and media (autoclave, laminar-flow
hood); observation of bacteria and the Gram
stain; bacteria and their environment
(sampling, serial dilution and plating);
selective media; biotypes; growth kinetics;
mutation, adaptation and repair (nalidixic
acid, which damages DNA and induces the SOS
response); genetic variability measured by
isoenzyme electrophoresis; and a phylogeny of
the isolates built by UPGMA.

\section{The Vibrio project}
\label{section-vibrio-project}

The practicals fed a team project, presented
in the course's prokaryote week, whose goal
was to find, isolate and characterize two
Gram-negative genera, \emph{Vibrio} and
\emph{Klebsiella}, from nature. The team
sampled stagnant water, soil, and faeces from
a dead frog in the Botanical Garden of UNAM's
Institute of Biology, grew the samples in a
propagation medium, plated serial dilutions,
and tested the colonies by Gram stain,
selective media (MacConkey, EMB, TCBS) and
biotype tests.

What the textbooks say about the two genera:
\begin{itemize}
\item \emph{Vibrio}: Gram-negative straight
or curved rods, mostly with a polar
flagellum, fermentative and oxidase-positive
(unlike the enteric bacteria), living in
marine, brackish and fresh water
\cite[Section 17.12]{madigan2012brock}.
\emph{V. cholerae} causes cholera, a watery
diarrhoea caused by a toxin and spread by
contaminated water; \emph{V.
parahaemolyticus} causes gastroenteritis from
raw seafood. The luminous squid symbiont that
the team mentioned as \emph{V. fischeri} is
now called \emph{Aliivibrio fischeri}; it
glows only at high cell density, through
quorum sensing \cite[Sections 8.9 and
17.12]{madigan2012brock}.
\item \emph{Klebsiella}: non-motile enteric
rods with a large polysaccharide capsule,
common in soil and water; many fix nitrogen,
and \emph{K. pneumoniae} occasionally causes
a serious pneumonia \cite[Chapter
11]{tortora2013microbiology}.
\end{itemize}
Result: the colonies grew well on the
selective and minimal media, and the
biochemical tests indicated \emph{Vibrio} but
excluded \emph{Klebsiella}, so the team could
not confirm that it had isolated both genera;
genetic tests were proposed as the next step.

\section{Notable genera}
\label{section-genera}

The final exam asked about genera presented
by students in class. The table below lists
them, with what the textbooks say about each.

\begin{center}
\begin{tabular}{p{0.3\linewidth}
p{0.6\linewidth}}
Genus & Why it matters \\
\hline
\emph{Rhizobium} & fixes $\mathrm{N_2}$ with
nitrogenase in root nodules of legumes
\cite[Section 25.3]{madigan2012brock} \\
\emph{Agrobacterium} & transfers DNA from its
Ti plasmid into plant cells (crown gall);
used to make transgenic plants \cite[Section
25.4]{madigan2012brock} \\
\emph{Vibrio cholerae} & cholera \\
\emph{Helicobacter pylori} & colonizes the
stomach mucosa (its urease neutralizes acid
locally); gastritis and peptic ulcers
\cite[Section 33.10]{madigan2012brock} \\
\emph{Staphylococcus aureus} & boils,
pneumonia, osteomyelitis and other infections
\cite[Chapter 18]{madigan2012brock} \\
\emph{Bacillus anthracis} & anthrax;
endospore former, a potential biological
weapon \cite[Chapter
11]{tortora2013microbiology} \\
\emph{Yersinia pestis} & plague, carried by
rodents and transmitted by fleas
\cite[Chapter 11]{tortora2013microbiology} \\
\emph{Erwinia}, \emph{Ralstonia} & plant
pathogens (soft rots; \emph{R. solanacearum})
\cite[Chapter 11]{tortora2013microbiology}
\cite[Section 17.7]{madigan2012brock} \\
\emph{Pseudomonas} (\emph{P. putida} in the
course) & degrades many organic pollutants,
including oil: bioremediation \cite[Chapter
2]{tortora2013microbiology} \\
\emph{Zoogloea} & forms the flocs of
activated-sludge sewage treatment
\cite[Chapter 11]{tortora2013microbiology} \\
\emph{Thiobacillus},
\emph{Acidithiobacillus} & sulfur oxidizers;
\emph{A. ferrooxidans} leaches copper ores
\cite[Section 24.7]{madigan2012brock} \\
\emph{Zymomonas} & ferments glucose to
ethanol by the Entner--Doudoroff pathway;
fermented plant saps \cite[Section
17.7]{madigan2012brock} \\
\emph{Bacteroides} & numerically dominant
bacteria of the human large intestine
\cite[Chapter 18]{madigan2012brock} \\
\emph{Wolbachia} & intracellular symbionts of
insects, ``probably the most common
infectious bacterial genus in the world''
\cite[Chapter 11]{tortora2013microbiology} \\
\emph{Pyrococcus} & archaeon, optimum 100
$^\circ$C \\
\emph{Methanosarcina} & archaeon, methanogen
\\
\emph{Halobacterium} & archaeon, extreme
halophile with bacteriorhodopsin \\
\end{tabular}
\end{center}

\section{Biogeochemical cycles and pollution}
\label{section-cycles}

Prokaryotes run most steps of the cycles of
the elements \cite[Chapter
24]{madigan2012brock}. In the nitrogen cycle,
nitrogen fixation turns $\mathrm{N_2}$ into
ammonia; nitrification oxidizes ammonium to
nitrite and nitrate (chemolithotrophic
Bacteria and Archaea); denitrification and
anammox return nitrogen to the air as
$\mathrm{N_2}$. In the carbon cycle,
methanogens make methane in anoxic habitats.
Sulfur and iron bacteria oxidize and reduce
sulfur and iron compounds; the oxidation of
iron sulfides by \emph{Acidithiobacillus
ferrooxidans} explains both microbial
leaching and acid mine drainage.

Daniel's essay for the course, based on a
review by Coelho and colleagues on marine
microbes, argued that climate change and
pollution act on ecosystems through these
cycles: the rise of atmospheric
$\mathrm{CO_2}$ acidifies water, ultraviolet
radiation and pollutants change toxicity, and
these changes alter the microbes that
transform nitrogen (nitrification depends on
the form of ammonium available) and iron.
Since bacteria are the main processors of
pollutants, and marine microbes were still
little studied, he concluded that
understanding them is necessary to predict
how ecosystems will change.

\bibliography{my}
\bibliographystyle{amsalpha}

\end{document}
